\chapter{Módulos de la plataforma}
\label{cap:modulos}

A lo largo de este capítulo se detallan los aspectos más relevantes de cada uno de los microservicios desplegados en la plataforma de ``Smart City'' de San Lorenzo de El Escorial.

\section{API Gateway (\texttt{atlas-api})}

Este microservicio forma parte de los servicios de ``vista'' (\ref{sec:microservicios}) y está compuesto por dos contenedores (\ref{cap:despliegue}): un contenedor basado en NGINX que actúa como proxy inverso y terminador SSL, y un contenedor de aplicación que implementa la API REST para aplicaciones ciudadanas y servicios municipales.

Se implementa en Node.js sobre Fastify 5 y expone su interfaz en el puerto \texttt{3000}. Las peticiones HTTP recibidas se resuelven mediante definición declarativa de endpoints y se traducen a mensajes AMQP usando \textit{ferimer-endpoints-manager} (\ref{sec:ferimer-endpoints-manager}) y \textit{ferimer-queue-manager} (\ref{sec:ferimer-queue-manager}).

Cada endpoint define su contrato de entrada/salida en esquemas JSON validados con AJV. Para operaciones síncronas se utiliza un emisor RPC con \textit{correlation id}, mientras que para eventos de notificación se emplea publicación \textit{Fire-and-Forget}. El resultado es una puerta de enlace HTTP desacoplada de la lógica de negocio y alineada con la ontología de mensajería (\ref{sec:atlas-messaging-ontology}).

\section{CMS (\texttt{atlas-cms})}

\texttt{atlas-cms} es el sistema de gestión de contenidos editorial y se basa en Strapi 5. Opera sobre MariaDB para persistencia relacional y utiliza Valkey como capa de caché para reducir latencia en lecturas frecuentes.

El servicio se publica en el puerto \texttt{1337} y soporta gestión multilenguaje (i18n), permitiendo mantener contenidos en varios idiomas desde un mismo panel. Además, el CMS puede incorporar fuentes externas de interés municipal (por ejemplo, información meteorológica de AEMET) para enriquecer páginas y componentes consumidos por la web pública.

Desde la perspectiva arquitectónica, el CMS se integra con el resto del ecosistema sin romper el desacoplamiento: las operaciones editoriales persisten en su plano de datos y las integraciones con otros módulos se realizan por API o mensajería según el caso de uso.

\section{Servicio de Inteligencia Artificial (\texttt{atlas-service-ollama})}

\texttt{atlas-service-ollama} centraliza las capacidades de inteligencia artificial de la plataforma y actúa como consumidor AMQP especializado. El servicio encapsula el acceso a Ollama para tres familias funcionales:

\begin{itemize}
\item \textbf{LLM}: generación y transformación de texto.
\item \textbf{OCR}: extracción de texto desde imágenes o documentos.
\item \textbf{Embeddings}: vectorización para búsqueda semántica y similitud.
\end{itemize}

Las peticiones llegan por colas RPC definidas por \textit{routing keys} derivadas de los tags de la ontología (por ejemplo, \texttt{tools:ai:llm} o \texttt{tools:ocr:image:file}). El servicio opera como \textit{RPC receiver}, valida payloads contra esquemas oficiales y devuelve respuestas estructuradas para mantener contratos estables entre productores y consumidores.

\section{Scraping de eventos (\texttt{atlas-calendar-scrapping})}

\texttt{atlas-calendar-scrapping} se encarga de capturar y normalizar información de eventos desde fuentes heterogéneas (sitios web, calendarios y redes sociales) para su posterior publicación en la plataforma.

La arquitectura del servicio sigue un patrón de especialización por adaptador:

\begin{itemize}
\item \texttt{BaseScrapper}: contrato común para extracción, transformación y normalización.
\item \texttt{ICSScrapper}: integración de calendarios en formato ICS.
\item \texttt{FBScrapper}: extracción de eventos desde fuentes basadas en Facebook.
\end{itemize}

Cada \textit{scrapper} implementa su lógica específica de autenticación, parseo y tratamiento de fechas, pero publica una salida uniforme para facilitar validación y consumo por API/CMS. Este enfoque permite añadir nuevas fuentes sin modificar el pipeline completo: basta con extender \texttt{BaseScrapper} y registrar el nuevo adaptador.

\section{Librerías de infraestructura}

\subsection{\texttt{ferimer-queue-manager}}\label{sec:ferimer-queue-manager}

(\url{https://github.com/Proyecto-Atlas/Ferimer-Queue-Manager})

Librería de abstracción AMQP utilizada por los microservicios de Atlas para estandarizar conexión, publicación y consumo de mensajes sobre RabbitMQ.

Expone cuatro clases base para cubrir los patrones de integración principales:

\begin{itemize}
\item \texttt{AMQP\_RPC\_Emitter}
\item \texttt{AMQP\_RPC\_Receiver}
\item \texttt{AMQP\_FireAndForget\_Emitter}
\item \texttt{AMQP\_FireAndForget\_Receiver}
\end{itemize}

La configuración habitual de Atlas utiliza el vhost \texttt{/atlas}, junto con nomenclatura de exchanges/colas alineada con los tags de la ontología de mensajes.

\subsection{\texttt{ferimer-endpoints-manager}}\label{sec:ferimer-endpoints-manager}

Librería para carga dinámica de endpoints HTTP en Fastify a partir de definiciones declarativas. Permite desacoplar la descripción de rutas de su implementación interna de mensajería.

Incluye validación de entrada y salida mediante esquemas JSON (AJV), soporte de parámetros de ruta (\textit{path params}) y mapeo de cada endpoint a su correspondiente tag/routing key AMQP.

Este mecanismo reduce código repetitivo en la API Gateway y facilita incorporar nuevas operaciones sin modificar el núcleo del servidor HTTP.

(\url{https://github.com/Proyecto-Atlas/Ferimer-Endpoints-Manager})

\subsection{\texttt{ferimer-database-manager}}

Librería de acceso a datos que encapsula la interacción con MariaDB y homogeniza el uso de conexión en los distintos servicios de Atlas.

Sus objetivos principales son:

\begin{itemize}
\item ofrecer una interfaz única para operaciones SQL,
\item minimizar errores de composición de consultas mediante utilidades de escape (\texttt{sqlstring}),
\item gestionar \textit{connection pooling} para mejorar rendimiento y estabilidad.
\end{itemize}

Esta capa simplifica el desarrollo de módulos de persistencia y centraliza buenas prácticas de acceso a base de datos.

\subsection{\texttt{atlas-messaging-ontology}}\label{sec:atlas-messaging-ontology}

(\url{https://github.com/Proyecto-Atlas/atlas-messaging-ontology})

Registro centralizado de esquemas JSON para los mensajes que circulan por las colas AMQP.
Define el formato canónico de etiquetas (\textit{tags}) con la estructura \texttt{dominio:categoría:acción[:subtipo]}
y proporciona validación AJV tanto para peticiones como para respuestas.

El catálogo de etiquetas del siguiente apartado se genera automáticamente
desde los esquemas JSON del repositorio fuente ejecutando:

\begin{lstlisting}[language=bash, caption={Regenerar catálogo de etiquetas}]
node scripts/generate-ontology-latex.js
\end{lstlisting}

% Catálogo generado automáticamente — NO editar ontologia-generada.tex a mano
% ─────────────────────────────────────────────────────────────────
% FICHERO GENERADO AUTOMÁTICAMENTE
% Comando: node scripts/generate-ontology-latex.js
% Fuente:  atlas-messaging-ontology/schemas/
% Etiquetas registradas: 5
% No editar manualmente — los cambios se perderán en la siguiente generación.
% ─────────────────────────────────────────────────────────────────

% Requiere en main.tex: \usepackage{longtable}, \usepackage{amssymb}

\subsection{Catálogo de etiquetas}
\label{sec:ontologia-catalogo}

El catálogo recoge las 5 etiquetas actualmente registradas en la ontología, organizadas por dominio funcional. Este contenido se genera automáticamente desde los esquemas JSON del repositorio \texttt{atlas-messaging-ontology}.

\subsubsection{Dominio \texttt{tools}}
\label{sec:ontologia-dominio-tools}

\subparagraph{\texttt{tools:ai:embedding}}
\label{sec:tag-tools-ai-embedding}

Dominio: \texttt{tools} — Categoría: \texttt{ai} — Acción: \texttt{embedding}

\paragraph{Petición (request)}

\textbf{Tag:} \texttt{tools:ai:embedding:request}\\
\textit{Request to generate vector embeddings from text input}

\begin{center}
\begin{longtable}{@{} l l c p{4.5cm} p{4cm} @{}}
\toprule
\textbf{Campo} & \textbf{Tipo} & \textbf{Req.} & \textbf{Descripción} & \textbf{Restricciones} \\
\midrule
\endfirsthead
\midrule
\textbf{Campo} & \textbf{Tipo} & \textbf{Req.} & \textbf{Descripción} & \textbf{Restricciones} \\
\midrule
\endhead
\midrule
\multicolumn{5}{r}{\footnotesize\itshape Continúa en la siguiente página} \\
\endfoot
\bottomrule
\endlastfoot
  \texttt{model} & \texttt{string} &  & Embedding model identifier & default: \texttt{embeddinggemma} \\
  \texttt{input} & \texttt{string / array} & \checkmark & Text or array of texts to generate embeddings for &  \\
  \texttt{dimensions} & \texttt{integer} &  & Desired dimensionality of the output vectors & mín: 64; máx: 4096 \\
  \texttt{correlation\_id} & \texttt{string} &  & Unique identifier to correlate request with response &  \\
  \texttt{reply\_to} & \texttt{string} &  & Queue tag where response should be sent & patrón: \texttt{\textasciicircum{}[a-z]+:[a-z]+:[a-z]+(:[a-z]+)?\$} \\
\end{longtable}
\end{center}

\paragraph{Respuesta (response)}

\textbf{Tag:} \texttt{tools:ai:embedding:response}\\
\textit{Response containing generated vector embeddings from an embedding model (e.g. Ollama)}

\begin{center}
\begin{longtable}{@{} l l c p{4.5cm} p{4cm} @{}}
\toprule
\textbf{Campo} & \textbf{Tipo} & \textbf{Req.} & \textbf{Descripción} & \textbf{Restricciones} \\
\midrule
\endfirsthead
\midrule
\textbf{Campo} & \textbf{Tipo} & \textbf{Req.} & \textbf{Descripción} & \textbf{Restricciones} \\
\midrule
\endhead
\midrule
\multicolumn{5}{r}{\footnotesize\itshape Continúa en la siguiente página} \\
\endfoot
\bottomrule
\endlastfoot
  \texttt{model} & \texttt{string} & \checkmark & Embedding model used for generation &  \\
  \texttt{embeddings} & \texttt{array[]} & \checkmark & Array of embedding vectors, one per input & minItems: 1 \\
  \texttt{total\_duration} & \texttt{integer} &  & Total time taken for the request in nanoseconds & mín: 0 \\
  \texttt{load\_duration} & \texttt{integer} &  & Time taken to load the model in nanoseconds & mín: 0 \\
  \texttt{prompt\_eval\_count} & \texttt{integer} &  & Number of tokens evaluated in the prompt & mín: 0 \\
  \texttt{embeddingPacked} & \texttt{string} &  & Base64-encoded packed representation of the embeddings &  \\
\end{longtable}
\end{center}


\subparagraph{\texttt{tools:ai:llm}}
\label{sec:tag-tools-ai-llm}

Dominio: \texttt{tools} — Categoría: \texttt{ai} — Acción: \texttt{llm}

\paragraph{Petición (request)}

\textbf{Tag:} \texttt{tools:ai:llm:request}\\
\textit{Request to generate text using a Large Language Model}

\begin{center}
\begin{longtable}{@{} l l c p{4.5cm} p{4cm} @{}}
\toprule
\textbf{Campo} & \textbf{Tipo} & \textbf{Req.} & \textbf{Descripción} & \textbf{Restricciones} \\
\midrule
\endfirsthead
\midrule
\textbf{Campo} & \textbf{Tipo} & \textbf{Req.} & \textbf{Descripción} & \textbf{Restricciones} \\
\midrule
\endhead
\midrule
\multicolumn{5}{r}{\footnotesize\itshape Continúa en la siguiente página} \\
\endfoot
\bottomrule
\endlastfoot
  \texttt{model} & \texttt{string} &  & LLM model identifier & default: \texttt{gemma3} \\
  \texttt{system\_prompt} & \texttt{string} &  & System instructions for the model &  \\
  \texttt{prompt} & \texttt{string} & \checkmark & User prompt to process &  \\
  \texttt{correlation\_id} & \texttt{string} &  & Unique identifier to correlate request with response &  \\
  \texttt{reply\_to} & \texttt{string} &  & Queue tag where response should be sent & patrón: \texttt{\textasciicircum{}[a-z]+:[a-z]+:[a-z]+(:[a-z]+)?\$} \\
\end{longtable}
\end{center}

\paragraph{Respuesta (response)}

\textbf{Tag:} \texttt{tools:ai:llm:response}\\
\textit{Response from a Large Language Model after processing a prompt}

\begin{center}
\begin{longtable}{@{} l l c p{4.5cm} p{4cm} @{}}
\toprule
\textbf{Campo} & \textbf{Tipo} & \textbf{Req.} & \textbf{Descripción} & \textbf{Restricciones} \\
\midrule
\endfirsthead
\midrule
\textbf{Campo} & \textbf{Tipo} & \textbf{Req.} & \textbf{Descripción} & \textbf{Restricciones} \\
\midrule
\endhead
\midrule
\multicolumn{5}{r}{\footnotesize\itshape Continúa en la siguiente página} \\
\endfoot
\bottomrule
\endlastfoot
  \texttt{model} & \texttt{string} & \checkmark & LLM model identifier used for generation &  \\
  \texttt{text} & \texttt{string} & \checkmark & Generated text response &  \\
  \texttt{usage} & \texttt{object} &  & Token usage statistics &  \\
  \texttt{finish\_reason} & \texttt{string} &  & Reason the model stopped generating & \texttt{stop}, \texttt{length}, \texttt{error} \\
  \texttt{processing\_time} & \texttt{integer} &  & Processing time &  \\
  \texttt{correlation\_id} & \texttt{string} &  & Identifier from the original request for correlation &  \\
\end{longtable}
\end{center}


\subparagraph{\texttt{tools:ai:translation}}
\label{sec:tag-tools-ai-translation}

Dominio: \texttt{tools} — Categoría: \texttt{ai} — Acción: \texttt{translation}

\paragraph{Petición (request)}

\textbf{Tag:} \texttt{tools:ai:translation:request}\\
\textit{Request to translate content to a target locale, typically triggered by a CMS i18n workflow}

\begin{center}
\begin{longtable}{@{} l l c p{4.5cm} p{4cm} @{}}
\toprule
\textbf{Campo} & \textbf{Tipo} & \textbf{Req.} & \textbf{Descripción} & \textbf{Restricciones} \\
\midrule
\endfirsthead
\midrule
\textbf{Campo} & \textbf{Tipo} & \textbf{Req.} & \textbf{Descripción} & \textbf{Restricciones} \\
\midrule
\endhead
\midrule
\multicolumn{5}{r}{\footnotesize\itshape Continúa en la siguiente página} \\
\endfoot
\bottomrule
\endlastfoot
  \texttt{jobId} & \texttt{string} & \checkmark & Unique job identifier composed of contentType and documentId (e.g., 'article:42') & patrón: \texttt{\textasciicircum{}[a-zA-Z0-9.\_-]+:[a-zA-Z0-9.\_-]+\$} \\
  \texttt{timestamp} & \texttt{integer} & \checkmark & Unix timestamp in milliseconds when the request was created & mín: 0 \\
  \texttt{contentType} & \texttt{string} & \checkmark & Content type identifier (e.g., Strapi UID like 'article', 'page') & minLen: 1 \\
  \texttt{documentId} & \texttt{string} & \checkmark & Unique identifier of the document to translate & minLen: 1 \\
  \texttt{targetLocale} & \texttt{string} & \checkmark & BCP 47 locale code for the target language (e.g., 'es', 'en-US', 'fr-FR') & patrón: \texttt{\textasciicircum{}[a-z]\{2\}(-[A-Z]\{2\})?\$} \\
  \texttt{data} & \texttt{object} & \checkmark & Key-value pairs of field names and their content to translate &  \\
  \texttt{prefix} & \texttt{string} &  & Prefix marker for pending translations (e.g., '[[TRANSLATION\_PENDING]]') & default: \texttt{[[TRANSLATION\_PENDING]]} \\
  \texttt{correlation\_id} & \texttt{string} &  & Unique identifier to correlate request with response &  \\
  \texttt{reply\_to} & \texttt{string} &  & Queue tag where response should be sent & patrón: \texttt{\textasciicircum{}[a-z]+:[a-z]+:[a-z]+(:[a-z]+)?\$} \\
\end{longtable}
\end{center}

\paragraph{Respuesta (response)}

\textbf{Tag:} \texttt{tools:ai:translation:response}\\
\textit{Response containing the translated content for a given document and target locale}

\begin{center}
\begin{longtable}{@{} l l c p{4.5cm} p{4cm} @{}}
\toprule
\textbf{Campo} & \textbf{Tipo} & \textbf{Req.} & \textbf{Descripción} & \textbf{Restricciones} \\
\midrule
\endfirsthead
\midrule
\textbf{Campo} & \textbf{Tipo} & \textbf{Req.} & \textbf{Descripción} & \textbf{Restricciones} \\
\midrule
\endhead
\midrule
\multicolumn{5}{r}{\footnotesize\itshape Continúa en la siguiente página} \\
\endfoot
\bottomrule
\endlastfoot
  \texttt{jobId} & \texttt{string} & \checkmark & Unique job identifier from the original request & patrón: \texttt{\textasciicircum{}[a-zA-Z0-9.\_-]+:[a-zA-Z0-9.\_-]+\$} \\
  \texttt{timestamp} & \texttt{integer} & \checkmark & Unix timestamp in milliseconds when the response was created & mín: 0 \\
  \texttt{contentType} & \texttt{string} & \checkmark & Content type identifier from the original request & minLen: 1 \\
  \texttt{documentId} & \texttt{string} & \checkmark & Unique identifier of the translated document & minLen: 1 \\
  \texttt{targetLocale} & \texttt{string} & \checkmark & BCP 47 locale code of the translated content & patrón: \texttt{\textasciicircum{}[a-z]\{2\}(-[A-Z]\{2\})?\$} \\
  \texttt{data} & \texttt{object} & \checkmark & Key-value pairs of field names and their translated content &  \\
  \texttt{status} & \texttt{string} & \checkmark & Result status of the translation job & \texttt{completed}, \texttt{partial}, \texttt{error} \\
  \texttt{error} & \texttt{object} &  & Error details when status is 'error' or 'partial' &  \\
  \texttt{correlation\_id} & \texttt{string} &  & Identifier from the original request for correlation &  \\
\end{longtable}
\end{center}


\subparagraph{\texttt{tools:ocr:image:file}}
\label{sec:tag-tools-ocr-image-file}

Dominio: \texttt{tools} — Categoría: \texttt{ocr} — Acción: \texttt{image:file}

\paragraph{Petición (request)}

\textbf{Tag:} \texttt{tools:ocr:image:file:request}\\
\textit{Request to extract text from an image sent as base64-encoded file}

\begin{center}
\begin{longtable}{@{} l l c p{4.5cm} p{4cm} @{}}
\toprule
\textbf{Campo} & \textbf{Tipo} & \textbf{Req.} & \textbf{Descripción} & \textbf{Restricciones} \\
\midrule
\endfirsthead
\midrule
\textbf{Campo} & \textbf{Tipo} & \textbf{Req.} & \textbf{Descripción} & \textbf{Restricciones} \\
\midrule
\endhead
\midrule
\multicolumn{5}{r}{\footnotesize\itshape Continúa en la siguiente página} \\
\endfoot
\bottomrule
\endlastfoot
  \texttt{content} & \texttt{string} & \checkmark & Base64-encoded image content &  \\
  \texttt{mimetype} & \texttt{string} & \checkmark & MIME type of the image & \texttt{image/png}, \texttt{image/jpeg}, \texttt{image/tiff}, \texttt{image/webp}, \texttt{image/bmp} \\
  \texttt{filename} & \texttt{string} &  & Original filename for reference &  \\
  \texttt{language} & \texttt{string} &  & Expected language of the text (ISO 639-1) & default: \texttt{es}; patrón: \texttt{\textasciicircum{}[a-z]\{2\}\$} \\
  \texttt{output\_format} & \texttt{string} &  & Desired output format for the extracted text & \texttt{text}, \texttt{json}, \texttt{hocr}; default: \texttt{text} \\
  \texttt{regions} & \texttt{object[]} &  & Specific regions of the image to process (coordinates in pixels) &  \\
  \texttt{correlation\_id} & \texttt{string} &  & Unique identifier to correlate request with response &  \\
  \texttt{reply\_to} & \texttt{string} &  & Queue tag where response should be sent & patrón: \texttt{\textasciicircum{}[a-z]+:[a-z]+:[a-z]+(:[a-z]+)?\$} \\
\end{longtable}
\end{center}

\paragraph{Respuesta (response)}

\textbf{Tag:} \texttt{tools:ocr:image:file:response}\\
\textit{Response containing text extracted from a base64-encoded image, as returned by an Ollama-compatible OCR model}

\begin{center}
\begin{longtable}{@{} l l c p{4.5cm} p{4cm} @{}}
\toprule
\textbf{Campo} & \textbf{Tipo} & \textbf{Req.} & \textbf{Descripción} & \textbf{Restricciones} \\
\midrule
\endfirsthead
\midrule
\textbf{Campo} & \textbf{Tipo} & \textbf{Req.} & \textbf{Descripción} & \textbf{Restricciones} \\
\midrule
\endhead
\midrule
\multicolumn{5}{r}{\footnotesize\itshape Continúa en la siguiente página} \\
\endfoot
\bottomrule
\endlastfoot
  \texttt{model} & \texttt{string} & \checkmark & Name of the OCR model used &  \\
  \texttt{created\_at} & \texttt{string} & \checkmark & Timestamp of when the response was generated (ISO 8601) & formato: \texttt{date-time} \\
  \texttt{response} & \texttt{string} & \checkmark & Extracted text content, may include markdown formatting &  \\
  \texttt{done} & \texttt{boolean} & \checkmark & Whether the generation is complete &  \\
  \texttt{done\_reason} & \texttt{string} & \checkmark & Reason why the generation stopped (e.g. 'stop') &  \\
  \texttt{context} & \texttt{integer[]} &  & Token context used during generation &  \\
  \texttt{total\_duration} & \texttt{integer} &  & Total processing time in nanoseconds & mín: 0 \\
  \texttt{load\_duration} & \texttt{integer} &  & Model load time in nanoseconds & mín: 0 \\
  \texttt{prompt\_eval\_count} & \texttt{integer} &  & Number of tokens in the prompt & mín: 0 \\
  \texttt{prompt\_eval\_duration} & \texttt{integer} &  & Time spent evaluating the prompt in nanoseconds & mín: 0 \\
  \texttt{eval\_count} & \texttt{integer} &  & Number of tokens generated in the response & mín: 0 \\
  \texttt{eval\_duration} & \texttt{integer} &  & Time spent generating the response in nanoseconds & mín: 0 \\
\end{longtable}
\end{center}


\subparagraph{\texttt{tools:ocr:image:url}}
\label{sec:tag-tools-ocr-image-url}

Dominio: \texttt{tools} — Categoría: \texttt{ocr} — Acción: \texttt{image:url}

\paragraph{Petición (request)}

\textbf{Tag:} \texttt{tools:ocr:image:url:request}\\
\textit{Request to extract text from an image referenced by URL}

\begin{center}
\begin{longtable}{@{} l l c p{4.5cm} p{4cm} @{}}
\toprule
\textbf{Campo} & \textbf{Tipo} & \textbf{Req.} & \textbf{Descripción} & \textbf{Restricciones} \\
\midrule
\endfirsthead
\midrule
\textbf{Campo} & \textbf{Tipo} & \textbf{Req.} & \textbf{Descripción} & \textbf{Restricciones} \\
\midrule
\endhead
\midrule
\multicolumn{5}{r}{\footnotesize\itshape Continúa en la siguiente página} \\
\endfoot
\bottomrule
\endlastfoot
  \texttt{url} & \texttt{string} & \checkmark & Public URL of the image to process & formato: \texttt{uri} \\
  \texttt{language} & \texttt{string} &  & Expected language of the text (ISO 639-1) & default: \texttt{es}; patrón: \texttt{\textasciicircum{}[a-z]\{2\}\$} \\
  \texttt{output\_format} & \texttt{string} &  & Desired output format for the extracted text & \texttt{text}, \texttt{json}, \texttt{hocr}; default: \texttt{text} \\
  \texttt{regions} & \texttt{object[]} &  & Specific regions of the image to process (coordinates in pixels) &  \\
  \texttt{correlation\_id} & \texttt{string} &  & Unique identifier to correlate request with response &  \\
  \texttt{reply\_to} & \texttt{string} &  & Queue tag where response should be sent & patrón: \texttt{\textasciicircum{}[a-z]+:[a-z]+:[a-z]+(:[a-z]+)?\$} \\
\end{longtable}
\end{center}

\paragraph{Respuesta (response)}

\textbf{Tag:} \texttt{tools:ocr:image:url:response}\\
\textit{Response containing text extracted from an image referenced by URL, as returned by an Ollama-compatible OCR model}

\begin{center}
\begin{longtable}{@{} l l c p{4.5cm} p{4cm} @{}}
\toprule
\textbf{Campo} & \textbf{Tipo} & \textbf{Req.} & \textbf{Descripción} & \textbf{Restricciones} \\
\midrule
\endfirsthead
\midrule
\textbf{Campo} & \textbf{Tipo} & \textbf{Req.} & \textbf{Descripción} & \textbf{Restricciones} \\
\midrule
\endhead
\midrule
\multicolumn{5}{r}{\footnotesize\itshape Continúa en la siguiente página} \\
\endfoot
\bottomrule
\endlastfoot
  \texttt{model} & \texttt{string} & \checkmark & Name of the OCR model used &  \\
  \texttt{created\_at} & \texttt{string} & \checkmark & Timestamp of when the response was generated (ISO 8601) & formato: \texttt{date-time} \\
  \texttt{response} & \texttt{string} & \checkmark & Extracted text content, may include markdown formatting &  \\
  \texttt{done} & \texttt{boolean} & \checkmark & Whether the generation is complete &  \\
  \texttt{done\_reason} & \texttt{string} & \checkmark & Reason why the generation stopped (e.g. 'stop') &  \\
  \texttt{context} & \texttt{integer[]} &  & Token context used during generation &  \\
  \texttt{total\_duration} & \texttt{integer} &  & Total processing time in nanoseconds & mín: 0 \\
  \texttt{load\_duration} & \texttt{integer} &  & Model load time in nanoseconds & mín: 0 \\
  \texttt{prompt\_eval\_count} & \texttt{integer} &  & Number of tokens in the prompt & mín: 0 \\
  \texttt{prompt\_eval\_duration} & \texttt{integer} &  & Time spent evaluating the prompt in nanoseconds & mín: 0 \\
  \texttt{eval\_count} & \texttt{integer} &  & Number of tokens generated in the response & mín: 0 \\
  \texttt{eval\_duration} & \texttt{integer} &  & Time spent generating the response in nanoseconds & mín: 0 \\
\end{longtable}
\end{center}


